% preamble.tex
% Hier werden alle LaTeX-Pakete geladen und konfiguriert.
% Diese Datei wird in main.tex mit % preamble.tex
% Hier werden alle LaTeX-Pakete geladen und konfiguriert.
% Diese Datei wird in main.tex mit % preamble.tex
% Hier werden alle LaTeX-Pakete geladen und konfiguriert.
% Diese Datei wird in main.tex mit % preamble.tex
% Hier werden alle LaTeX-Pakete geladen und konfiguriert.
% Diese Datei wird in main.tex mit \input{preamble} eingebunden.

% --- Schrift & Encoding ---
\usepackage[T1]{fontenc}       % Korrekte Silbentrennung und Sonderzeichen
\usepackage[utf8]{inputenc}    % UTF-8 Eingabe (Umlaute direkt schreiben)
\usepackage[ngerman]{babel}    % Deutsche Spracheinstellungen

% --- Seitenformat ---
\usepackage[a4paper,margin=2.5cm]{geometry}  % A4, 2,5 cm Rand rundum

% --- Zeilenabstand ---
\usepackage{setspace}
\onehalfspacing                % 1,5-facher Zeilenabstand (Nordakademie-Vorgabe)

% --- Kopf- und Fusszeilen ---
\usepackage{fancyhdr}
\setlength{\headheight}{15pt}
\fancyhf{}
\fancyhead[R]{\small\nouppercase{\leftmark}}  % Kapitelname oben rechts
\fancyfoot[C]{\thepage}                        % Seitenzahl unten mittig
\renewcommand{\headrulewidth}{0.4pt}

% --- Grafiken & Tabellen ---
\usepackage{graphicx}   % Bilder einbinden (\includegraphics)
\usepackage{pdfpages}   % PDF-Seiten einbinden (\includepdf) -- fuer Anhang
\usepackage{booktabs}   % Professionelle Tabellen (\toprule, \midrule, \bottomrule)
\usepackage{longtable}  % Tabellen ueber mehrere Seiten

% --- Quellcode ---
\usepackage{listings}
\lstset{
  basicstyle=\small\ttfamily,
  keywordstyle=\bfseries\color{blue!70!black},
  commentstyle=\itshape\color{gray},
  stringstyle=\color{green!50!black},
  breaklines=true,
  frame=single,
  numbers=left,
  numberstyle=\tiny\color{gray},
  captionpos=b,
  tabsize=2,
  showstringspaces=false,
  % Umlaute in Listings:
  literate={ü}{{\"{u}}}1 {ö}{{\"{o}}}1 {ä}{{\"{a}}}1
           {Ü}{{\"{U}}}1 {Ö}{{\"{O}}}1 {Ä}{{\"{A}}}1 {ß}{{\ss{}}}1
}

% --- Sonstiges ---
\usepackage{csquotes}   % Korrekte Anfuehrungszeichen (benoetigt von biblatex)
\usepackage{xcolor}     % Farben (z.B. fuer Links und Listings)
\usepackage{amssymb}    % Mathematische Symbole (\square fuer Checkboxen im Fragebogen)
\usepackage{acronym}    % Abkuerzungsverzeichnis (\ac{}, \acl{}, \acs{})

% --- Hyperlinks & PDF-Metadaten ---
\usepackage{hyperref}
\hypersetup{
  colorlinks=true,
  linkcolor=blue,    % Interne Links (Inhaltsverzeichnis etc.)
  citecolor=blue,    % Zitierlinks
  urlcolor=blue,     % URLs
  % !! Bitte anpassen !!
  pdftitle={Bachelorarbeit -- Titel der Arbeit},
  pdfauthor={Vorname Nachname},
  pdfsubject={Bachelorarbeit},
  pdfkeywords={Bachelorarbeit, Stichwort1, Stichwort2}
}

% --- Literaturverzeichnis ---
% biblatex mit biber als Backend, Autor-Jahr-Stil, Zitate als Fussnoten
\usepackage[
  backend=biber,        % Biber statt BibTeX (moderner, Umlaute-sicher)
  style=authoryear,     % Autor-Jahr-Zitierweise
  autocite=footnote,    % \autocite{} erzeugt automatisch Fussnoten
  isbn=false,           % ISBN nicht anzeigen
  doi=false,            % DOI nicht anzeigen
  url=true,             % URLs anzeigen
  giveninits=true,      % Vornamen abkuerzen (z.B. "M. Fowler")
  uniquename=init,
  maxcitenames=2,       % Im Text max. 2 Autoren, dann "u. a."
  maxbibnames=99,       % Im Literaturverzeichnis alle Autoren
  ibidtracker=false,
  pagetracker=false
]{biblatex}
\addbibresource{references.bib}  % Literatur-Datei einbinden
 eingebunden.

% --- Schrift & Encoding ---
\usepackage[T1]{fontenc}       % Korrekte Silbentrennung und Sonderzeichen
\usepackage[utf8]{inputenc}    % UTF-8 Eingabe (Umlaute direkt schreiben)
\usepackage[ngerman]{babel}    % Deutsche Spracheinstellungen

% --- Seitenformat ---
\usepackage[a4paper,margin=2.5cm]{geometry}  % A4, 2,5 cm Rand rundum

% --- Zeilenabstand ---
\usepackage{setspace}
\onehalfspacing                % 1,5-facher Zeilenabstand (Nordakademie-Vorgabe)

% --- Kopf- und Fusszeilen ---
\usepackage{fancyhdr}
\setlength{\headheight}{15pt}
\fancyhf{}
\fancyhead[R]{\small\nouppercase{\leftmark}}  % Kapitelname oben rechts
\fancyfoot[C]{\thepage}                        % Seitenzahl unten mittig
\renewcommand{\headrulewidth}{0.4pt}

% --- Grafiken & Tabellen ---
\usepackage{graphicx}   % Bilder einbinden (\includegraphics)
\usepackage{pdfpages}   % PDF-Seiten einbinden (\includepdf) -- fuer Anhang
\usepackage{booktabs}   % Professionelle Tabellen (\toprule, \midrule, \bottomrule)
\usepackage{longtable}  % Tabellen ueber mehrere Seiten

% --- Quellcode ---
\usepackage{listings}
\lstset{
  basicstyle=\small\ttfamily,
  keywordstyle=\bfseries\color{blue!70!black},
  commentstyle=\itshape\color{gray},
  stringstyle=\color{green!50!black},
  breaklines=true,
  frame=single,
  numbers=left,
  numberstyle=\tiny\color{gray},
  captionpos=b,
  tabsize=2,
  showstringspaces=false,
  % Umlaute in Listings:
  literate={ü}{{\"{u}}}1 {ö}{{\"{o}}}1 {ä}{{\"{a}}}1
           {Ü}{{\"{U}}}1 {Ö}{{\"{O}}}1 {Ä}{{\"{A}}}1 {ß}{{\ss{}}}1
}

% --- Sonstiges ---
\usepackage{csquotes}   % Korrekte Anfuehrungszeichen (benoetigt von biblatex)
\usepackage{xcolor}     % Farben (z.B. fuer Links und Listings)
\usepackage{amssymb}    % Mathematische Symbole (\square fuer Checkboxen im Fragebogen)
\usepackage{acronym}    % Abkuerzungsverzeichnis (\ac{}, \acl{}, \acs{})

% --- Hyperlinks & PDF-Metadaten ---
\usepackage{hyperref}
\hypersetup{
  colorlinks=true,
  linkcolor=blue,    % Interne Links (Inhaltsverzeichnis etc.)
  citecolor=blue,    % Zitierlinks
  urlcolor=blue,     % URLs
  % !! Bitte anpassen !!
  pdftitle={Bachelorarbeit -- Titel der Arbeit},
  pdfauthor={Vorname Nachname},
  pdfsubject={Bachelorarbeit},
  pdfkeywords={Bachelorarbeit, Stichwort1, Stichwort2}
}

% --- Literaturverzeichnis ---
% biblatex mit biber als Backend, Autor-Jahr-Stil, Zitate als Fussnoten
\usepackage[
  backend=biber,        % Biber statt BibTeX (moderner, Umlaute-sicher)
  style=authoryear,     % Autor-Jahr-Zitierweise
  autocite=footnote,    % \autocite{} erzeugt automatisch Fussnoten
  isbn=false,           % ISBN nicht anzeigen
  doi=false,            % DOI nicht anzeigen
  url=true,             % URLs anzeigen
  giveninits=true,      % Vornamen abkuerzen (z.B. "M. Fowler")
  uniquename=init,
  maxcitenames=2,       % Im Text max. 2 Autoren, dann "u. a."
  maxbibnames=99,       % Im Literaturverzeichnis alle Autoren
  ibidtracker=false,
  pagetracker=false
]{biblatex}
\addbibresource{references.bib}  % Literatur-Datei einbinden
 eingebunden.

% --- Schrift & Encoding ---
\usepackage[T1]{fontenc}       % Korrekte Silbentrennung und Sonderzeichen
\usepackage[utf8]{inputenc}    % UTF-8 Eingabe (Umlaute direkt schreiben)
\usepackage[ngerman]{babel}    % Deutsche Spracheinstellungen

% --- Seitenformat ---
\usepackage[a4paper,margin=2.5cm]{geometry}  % A4, 2,5 cm Rand rundum

% --- Zeilenabstand ---
\usepackage{setspace}
\onehalfspacing                % 1,5-facher Zeilenabstand (Nordakademie-Vorgabe)

% --- Kopf- und Fusszeilen ---
\usepackage{fancyhdr}
\setlength{\headheight}{15pt}
\fancyhf{}
\fancyhead[R]{\small\nouppercase{\leftmark}}  % Kapitelname oben rechts
\fancyfoot[C]{\thepage}                        % Seitenzahl unten mittig
\renewcommand{\headrulewidth}{0.4pt}

% --- Grafiken & Tabellen ---
\usepackage{graphicx}   % Bilder einbinden (\includegraphics)
\usepackage{pdfpages}   % PDF-Seiten einbinden (\includepdf) -- fuer Anhang
\usepackage{booktabs}   % Professionelle Tabellen (\toprule, \midrule, \bottomrule)
\usepackage{longtable}  % Tabellen ueber mehrere Seiten

% --- Quellcode ---
\usepackage{listings}
\lstset{
  basicstyle=\small\ttfamily,
  keywordstyle=\bfseries\color{blue!70!black},
  commentstyle=\itshape\color{gray},
  stringstyle=\color{green!50!black},
  breaklines=true,
  frame=single,
  numbers=left,
  numberstyle=\tiny\color{gray},
  captionpos=b,
  tabsize=2,
  showstringspaces=false,
  % Umlaute in Listings:
  literate={ü}{{\"{u}}}1 {ö}{{\"{o}}}1 {ä}{{\"{a}}}1
           {Ü}{{\"{U}}}1 {Ö}{{\"{O}}}1 {Ä}{{\"{A}}}1 {ß}{{\ss{}}}1
}

% --- Sonstiges ---
\usepackage{csquotes}   % Korrekte Anfuehrungszeichen (benoetigt von biblatex)
\usepackage{xcolor}     % Farben (z.B. fuer Links und Listings)
\usepackage{amssymb}    % Mathematische Symbole (\square fuer Checkboxen im Fragebogen)
\usepackage{acronym}    % Abkuerzungsverzeichnis (\ac{}, \acl{}, \acs{})

% --- Hyperlinks & PDF-Metadaten ---
\usepackage{hyperref}
\hypersetup{
  colorlinks=true,
  linkcolor=blue,    % Interne Links (Inhaltsverzeichnis etc.)
  citecolor=blue,    % Zitierlinks
  urlcolor=blue,     % URLs
  % !! Bitte anpassen !!
  pdftitle={Bachelorarbeit -- Titel der Arbeit},
  pdfauthor={Vorname Nachname},
  pdfsubject={Bachelorarbeit},
  pdfkeywords={Bachelorarbeit, Stichwort1, Stichwort2}
}

% --- Literaturverzeichnis ---
% biblatex mit biber als Backend, Autor-Jahr-Stil, Zitate als Fussnoten
\usepackage[
  backend=biber,        % Biber statt BibTeX (moderner, Umlaute-sicher)
  style=authoryear,     % Autor-Jahr-Zitierweise
  autocite=footnote,    % \autocite{} erzeugt automatisch Fussnoten
  isbn=false,           % ISBN nicht anzeigen
  doi=false,            % DOI nicht anzeigen
  url=true,             % URLs anzeigen
  giveninits=true,      % Vornamen abkuerzen (z.B. "M. Fowler")
  uniquename=init,
  maxcitenames=2,       % Im Text max. 2 Autoren, dann "u. a."
  maxbibnames=99,       % Im Literaturverzeichnis alle Autoren
  ibidtracker=false,
  pagetracker=false
]{biblatex}
\addbibresource{references.bib}  % Literatur-Datei einbinden
 eingebunden.

% --- Schrift & Encoding ---
\usepackage[T1]{fontenc}       % Korrekte Silbentrennung und Sonderzeichen
\usepackage[utf8]{inputenc}    % UTF-8 Eingabe (Umlaute direkt schreiben)
\usepackage[ngerman]{babel}    % Deutsche Spracheinstellungen

% --- Seitenformat ---
\usepackage[a4paper,margin=2.5cm]{geometry}  % A4, 2,5 cm Rand rundum

% --- Zeilenabstand ---
\usepackage{setspace}
\onehalfspacing                % 1,5-facher Zeilenabstand (Nordakademie-Vorgabe)

% --- Kopf- und Fusszeilen ---
\usepackage{fancyhdr}
\setlength{\headheight}{15pt}
\fancyhf{}
\fancyhead[R]{\small\nouppercase{\leftmark}}  % Kapitelname oben rechts
\fancyfoot[C]{\thepage}                        % Seitenzahl unten mittig
\renewcommand{\headrulewidth}{0.4pt}

% --- Grafiken & Tabellen ---
\usepackage{graphicx}   % Bilder einbinden (\includegraphics)
\usepackage{pdfpages}   % PDF-Seiten einbinden (\includepdf) -- fuer Anhang
\usepackage{booktabs}   % Professionelle Tabellen (\toprule, \midrule, \bottomrule)
\usepackage{longtable}  % Tabellen ueber mehrere Seiten

% --- Quellcode ---
\usepackage{listings}
\lstset{
  basicstyle=\small\ttfamily,
  keywordstyle=\bfseries\color{blue!70!black},
  commentstyle=\itshape\color{gray},
  stringstyle=\color{green!50!black},
  breaklines=true,
  frame=single,
  numbers=left,
  numberstyle=\tiny\color{gray},
  captionpos=b,
  tabsize=2,
  showstringspaces=false,
  % Umlaute in Listings:
  literate={ü}{{\"{u}}}1 {ö}{{\"{o}}}1 {ä}{{\"{a}}}1
           {Ü}{{\"{U}}}1 {Ö}{{\"{O}}}1 {Ä}{{\"{A}}}1 {ß}{{\ss{}}}1
}

% --- Sonstiges ---
\usepackage{csquotes}   % Korrekte Anfuehrungszeichen (benoetigt von biblatex)
\usepackage{xcolor}     % Farben (z.B. fuer Links und Listings)
\usepackage{amssymb}    % Mathematische Symbole (\square fuer Checkboxen im Fragebogen)
\usepackage{acronym}    % Abkuerzungsverzeichnis (\ac{}, \acl{}, \acs{})

% --- Hyperlinks & PDF-Metadaten ---
\usepackage{hyperref}
\hypersetup{
  colorlinks=true,
  linkcolor=blue,    % Interne Links (Inhaltsverzeichnis etc.)
  citecolor=blue,    % Zitierlinks
  urlcolor=blue,     % URLs
  % !! Bitte anpassen !!
  pdftitle={Bachelorarbeit -- Titel der Arbeit},
  pdfauthor={Vorname Nachname},
  pdfsubject={Bachelorarbeit},
  pdfkeywords={Bachelorarbeit, Stichwort1, Stichwort2}
}

% --- Literaturverzeichnis ---
% biblatex mit biber als Backend, Autor-Jahr-Stil, Zitate als Fussnoten
\usepackage[
  backend=biber,        % Biber statt BibTeX (moderner, Umlaute-sicher)
  style=authoryear,     % Autor-Jahr-Zitierweise
  autocite=footnote,    % \autocite{} erzeugt automatisch Fussnoten
  isbn=false,           % ISBN nicht anzeigen
  doi=false,            % DOI nicht anzeigen
  url=true,             % URLs anzeigen
  giveninits=true,      % Vornamen abkuerzen (z.B. "M. Fowler")
  uniquename=init,
  maxcitenames=2,       % Im Text max. 2 Autoren, dann "u. a."
  maxbibnames=99,       % Im Literaturverzeichnis alle Autoren
  ibidtracker=false,
  pagetracker=false
]{biblatex}
\addbibresource{references.bib}  % Literatur-Datei einbinden
